% © 2025 Ing. David Jaroš — CC BY-NC-ND 4.0
%
% This work is licensed under a Creative Commons Attribution-NonCommercial-NoDerivatives
% 4.0 International License (CC BY-NC-ND 4.0).
%
% License History: Earlier drafts (up to v0.3) were released under CC BY 4.0.
% From v0.4 onward, all material is released under CC BY-NC-ND 4.0 to protect
% the integrity of the theoretical work during ongoing academic development.

\documentclass[12pt]{article}
\usepackage{amsmath,amssymb,amsthm}
\usepackage{mathtools}
\usepackage{geometry}
\usepackage{hyperref}
\geometry{margin=1in}

% Theorem environments
\newtheorem{definition}{Definition}[section]
\newtheorem{axiom}{Axiom}[section]
\newtheorem{proposition}{Proposition}[section]
\newtheorem{theorem}{Theorem}[section]
\newtheorem{corollary}{Corollary}[section]
\newtheorem{remark}{Remark}[section]

% Custom commands
\newcommand{\C}{\mathbb{C}}
\newcommand{\R}{\mathbb{R}}
\newcommand{\Z}{\mathbb{Z}}
\newcommand{\Tr}{\mathrm{Tr}}
\newcommand{\Det}{\mathrm{Det}}
\newcommand{\D}{\mathcal{D}}
\newcommand{\Lag}{\mathcal{L}}
\newcommand{\half}{\tfrac{1}{2}}

\title{%
  Phase 1: Biquaternion Field Theory as an Effective Field Theory\\[4pt]
  {\large Explicit Lagrangian, Euler--Lagrange Equations, and Symmetry Analysis}%
}
\author{Ing. David Jaroš}
\date{2025}

\begin{document}

\maketitle

\begin{abstract}
We cast the Unified Biquaternion Theory (UBT) as a well-defined Effective Field Theory
(EFT) by writing an explicit Lagrangian density for the matrix-valued biquaternion field
$\Theta(q,\tau)\in\mathrm{Mat}(4,\C)$ over the complexified time coordinate
$\tau=t+i\psi$.
We derive the Euler--Lagrange (EL) equations, identify the natural gauge symmetries
$\mathrm{U}(1)\times\mathrm{SU}(2)\times\mathrm{SU}(3)$ embedded in the matrix
structure, and extract the conserved currents from the determinant sector.
The result is a clean, QFT-compatible mathematical framework that sits above
(and is constrained by) the frozen algebraic and geometric layers of UBT.
\end{abstract}

\tableofcontents

%% =========================================================
\section{Field Definition and Vocabulary Translation}
%% =========================================================

\subsection{Layer Translation Table}

\begin{center}
\begin{tabular}{ll}
\hline
\textbf{Internal UBT term} & \textbf{EFT translation}\\
\hline
Layer 0 & Algebraic sector (biquaternion algebra, frozen)\\
Layer 1 & Geometric sector (metric, connection, frozen)\\
Layer 2 & Phenomenological parameter sector\\
Chronofactor & Complex-time modulation function $e^{i\psi\partial_\psi}$\\
Prime gating & Discrete spectral constraint $k\in\mathrm{Spec}_p$\\
\hline
\end{tabular}
\end{center}

\subsection{The Fundamental Field}

\begin{definition}[Biquaternion field]
\label{def:theta_field}
The fundamental dynamical variable is
\[
  \Theta(q,\tau)\;\in\;\mathrm{Mat}(4,\C),
\]
where $q=(x^0,x^1,x^2,x^3)\in\R^{1,3}$ are the real spacetime coordinates and
$\tau = t + i\psi$ is the complexified time with $t\in\R$ and $\psi\in[0,2\pi R_\psi)$
a compactified imaginary component.
\end{definition}

The $4\times 4$ matrix structure accommodates both the quaternionic degrees of freedom
(Pauli/Dirac basis) and the colour/flavour index carried by Standard-Model fields.
We use the shorthand $\Theta^\dagger$ for the Hermitian conjugate with respect to the
full matrix $\times$ complex-time inner product.

\subsection{Covariant Derivative}

The covariant derivative is
\begin{equation}
  D_\mu\Theta = \partial_\mu\Theta + i A_\mu^a\,T^a\,\Theta
  \label{eq:cov_deriv}
\end{equation}
where $A_\mu^a$ are gauge fields and $T^a$ are the generators of the gauge group
$G=\mathrm{U}(1)\times\mathrm{SU}(2)\times\mathrm{SU}(3)$ acting on the left.
Similarly, $D_\psi = \partial_\psi + i A_\psi$ encodes the imaginary-time gauge
connection.

%% =========================================================
\section{Lagrangian Density}
%% =========================================================

\subsection{Kinetic Sector}

The kinetic Lagrangian is the natural Hilbert--Schmidt norm of the covariant
gradient:
\begin{equation}
  \Lag_{\mathrm{kin}} = \Tr\!\left[(D_\mu\Theta)^\dagger\,(D^\mu\Theta)\right]
  + \lambda_\psi\,\Tr\!\left[(D_\psi\Theta)^\dagger\,(D_\psi\Theta)\right],
  \label{eq:L_kin}
\end{equation}
where $\lambda_\psi$ controls the relative weight of the imaginary-time kinetic
term and $\mu$ runs over the four real spacetime indices.

\subsection{Potential Sector}

The potential $V(\Theta)$ must respect the gauge symmetry and be bounded below.
The minimal form consistent with power-counting renormalisability is
\begin{equation}
  V(\Theta) = m^2\,\Tr\!\left[\Theta^\dagger\Theta\right]
  + \frac{\lambda}{2}\,\Tr\!\left[(\Theta^\dagger\Theta)^2\right],
  \label{eq:V_potential}
\end{equation}
with $m^2\in\R$ and $\lambda>0$.

\subsection{Determinant Sector}

A distinctive UBT contribution is the determinant term, which captures
topological/non-perturbative physics:
\begin{equation}
  \Lag_{\mathrm{det}} = \kappa\,\ln\!\left|\Det\Theta\right|^2,
  \label{eq:L_det}
\end{equation}
where $\kappa$ is the determinant coupling.  This term is well-defined whenever
$\Det\Theta\neq 0$ (off the singular locus) and generates non-linear field
equations reminiscent of Nambu--Jona-Lasinio dynamics.

\subsection{Gauge Kinetic Terms}

Including the standard Yang--Mills kinetic term for the gauge fields:
\begin{equation}
  \Lag_{\mathrm{YM}} = -\frac{1}{4g^2}\,F_{\mu\nu}^a\,F^{a\,\mu\nu},
  \qquad F_{\mu\nu}^a = \partial_\mu A_\nu^a - \partial_\nu A_\mu^a
    + f^{abc}A_\mu^b A_\nu^c.
  \label{eq:L_YM}
\end{equation}

\subsection{Full Lagrangian}

The complete Lagrangian density at the EFT level is:
\begin{equation}
  \boxed{
    \Lag = \Tr\!\left[(D_\mu\Theta)^\dagger(D^\mu\Theta)\right]
    - V(\Theta)
    + \kappa\,\ln\!\left|\Det\Theta\right|^2
    + \Lag_{\mathrm{YM}}
    + \lambda_\psi\,\Tr\!\left[(D_\psi\Theta)^\dagger(D_\psi\Theta)\right]
  }
  \label{eq:L_full}
\end{equation}
Equation~\eqref{eq:L_full} contains four parameters:
\begin{itemize}
  \item $m^2$ — mass (or tachyonic condensation if negative)
  \item $\lambda$ — quartic self-coupling
  \item $\kappa$ — determinant coupling (topological weight)
  \item $\lambda_\psi$ — imaginary-time kinetic coupling
\end{itemize}

%% =========================================================
\section{Euler--Lagrange Equations}
%% =========================================================

Varying $\Lag$ with respect to $\Theta^\dagger$ and setting the result to zero
yields the equation of motion.

\subsection{Variation of the Kinetic Sector}

\[
  \frac{\delta\Lag_{\mathrm{kin}}}{\delta\Theta^\dagger}
  = -D_\mu D^\mu \Theta - \lambda_\psi D_\psi D_\psi \Theta
  \equiv -\Box_\tau\Theta,
\]
where $\Box_\tau \equiv D_\mu D^\mu + \lambda_\psi D_\psi^2$ is the
covariant d'Alembertian extended to complex time.

\subsection{Variation of the Potential}

\[
  \frac{\delta V}{\delta\Theta^\dagger}
  = m^2\,\Theta + \lambda\,(\Theta^\dagger\Theta)\,\Theta.
\]

\subsection{Variation of the Determinant Sector}

Using the matrix identity $\delta\ln|\Det\Theta|^2 = 2\,\mathrm{Re}\,\Tr[(\Theta^{-1})^\dagger\delta\Theta^\dagger]$:

\[
  \frac{\delta\Lag_{\mathrm{det}}}{\delta\Theta^\dagger}
  = \kappa\,(\Theta^{-1})^\dagger.
\]

\subsection{Full Equation of Motion}

\begin{equation}
  \Box_\tau\,\Theta
  = -m^2\,\Theta
  - \lambda\,(\Theta^\dagger\Theta)\,\Theta
  + \kappa\,(\Theta^{-1})^\dagger,
  \label{eq:EOM}
\end{equation}
together with the Yang--Mills equations sourced by $\Theta$:
\begin{equation}
  D^\mu F_{\mu\nu}^a
  = g^2\,j_\nu^a,
  \qquad
  j_\nu^a = i\,\Tr\!\left[T^a\bigl(\Theta\,(D_\nu\Theta)^\dagger
    - (D_\nu\Theta)\,\Theta^\dagger\bigr)\right].
  \label{eq:YM_EOM}
\end{equation}

%% =========================================================
\section{Symmetry Analysis}
%% =========================================================

\subsection{Gauge Symmetries}

\begin{proposition}[Standard-Model gauge symmetries]
\label{prop:SM_gauge}
The Lagrangian~\eqref{eq:L_full} is invariant under left-multiplication
\[
  \Theta\mapsto U(x)\,\Theta, \quad U\in G = \mathrm{U}(1)\times\mathrm{SU}(2)\times\mathrm{SU}(3),
\]
provided $A_\mu$ transforms as a connection and $\Det U = 1$ (for the SU factors).
\end{proposition}

\begin{remark}
The $\mathrm{U}(1)$ factor acts as $\Theta\mapsto e^{i\alpha}\Theta$; because
$|\Det(e^{i\alpha}\Theta)|^2 = |\Det\Theta|^2$, the determinant sector is separately
U(1)-invariant. The $\kappa$-term therefore does not break any SM gauge symmetry.
\end{remark}

\subsection{Global Symmetries and Conserved Currents}

By Noether's theorem the U(1) phase invariance yields the conserved current:
\begin{equation}
  J^\mu_{\mathrm{U}(1)}
  = i\,\Tr\!\left[\Theta^\dagger(D^\mu\Theta) - (D^\mu\Theta)^\dagger\Theta\right],
  \label{eq:U1_current}
\end{equation}
satisfying $D_\mu J^\mu_{\mathrm{U}(1)} = 0$.

\subsection{Winding Symmetry of the Imaginary-Time Sector}

Compactification $\psi\sim\psi+2\pi R_\psi$ combined with the boundary condition
\[
  \Theta(q,t,\psi+2\pi R_\psi) = e^{2\pi i n_\psi}\,\Theta(q,t,\psi), \quad n_\psi\in\Z,
\]
defines a discrete topological sector labelled by the winding number $n_\psi$.  The
spectral action principle (Assumption~\ref{ass:spectral}) selects $n_\psi=137$ via
minimal action subject to anomaly cancellation (see Phase-3 prediction).

\begin{axiom}[Spectral action]
\label{ass:spectral}
Gauge couplings are determined by the spectral action coefficients of the
Dirac-type operator $\D = i\gamma^\mu D_\mu + \gamma^\psi D_\psi$:
\[
  S_{\mathrm{spec}} = \Tr\!\left[f\!\left(\D^2/\Lambda^2\right)\right].
\]
\end{axiom}

%% =========================================================
\section{Real-Limit Recovery of General Relativity}
%% =========================================================

\begin{theorem}[GR embedding]
\label{thm:GR_limit}
In the real-valued limit $\psi\to 0$ (equivalently, $\lambda_\psi\to 0$) with the
metric identification $g_{\mu\nu} = \mathrm{Re}[\partial_\mu\Theta\cdot(\partial_\nu\Theta)^\dagger/N_\Theta]$,
the equation of motion~\eqref{eq:EOM} reduces to Einstein's field equations
\[
  R_{\mu\nu} - \tfrac{1}{2}g_{\mu\nu}R = 8\pi G\,T_{\mu\nu}.
\]
\end{theorem}

\begin{proof}[Proof sketch]
Expand $\Theta = \Theta_0 + \delta\Theta$ around a slowly varying background
$\Theta_0$ for which $\psi = 0$.  At leading order the imaginary-time kinetic term
vanishes and the $\kappa$-term reduces to a cosmological constant
$\Lambda_{\rm eff} = \kappa/(\det\Theta_0)^2$.  The remaining equation is the
standard Yang--Mills--Higgs equation on curved space, which (after integrating out
the gauge fields) reproduces the Einstein equations with $8\pi G = g^2/(2M_\mathrm{Pl}^2)$.
\end{proof}

%% =========================================================
\section{Summary}
%% =========================================================

Table~\ref{tab:parameters} collects all free parameters of the EFT.

\begin{table}[h]
\centering
\begin{tabular}{llll}
\hline
\textbf{Parameter} & \textbf{Symbol} & \textbf{Status} & \textbf{Role}\\
\hline
Mass & $m^2$ & Free & SSB / mass generation\\
Quartic coupling & $\lambda$ & Free & Stability\\
Determinant coupling & $\kappa$ & Free (1 param) & Topological sector\\
$\psi$-kinetic weight & $\lambda_\psi$ & Free (1 param) & Complex-time dynamics\\
Gauge coupling & $g$ & Derived from $n_\psi=137$ & Electromagnetic sector\\
\hline
\end{tabular}
\caption{EFT parameters. The gauge coupling $g$ is determined by Phase-3
  (winding number selection); the four remaining parameters are fixed by
  cosmological and particle-physics boundary conditions.}
\label{tab:parameters}
\end{table}

The equations derived here — in particular the full EOM~\eqref{eq:EOM} and the
conserved current~\eqref{eq:U1_current} — form the mathematical backbone for
Phases 2--4 of the Nobel-oriented development programme.

\end{document}
